\section{Deferred Proofs}\label{app:deferred-proofs}

\subsection{Unconditional proof of the rank-three lower bound}\label{app:t2-proof}

This appendix proves the unconditional lower bound in \Cref{thm:t2}.  The proof
is self-contained: all target families, auxiliary games, weights, activation
rules, deletion errors, and asymptotic estimates are defined below.  The
opening counterexamples diagnose why an older, stronger auxiliary potential
cannot prove the result.  The proof then uses potentials that count exact
target moves by either player, exactly as the game-length objective does.

Throughout this appendix $n$ is large, $0<\delta<1/4$ is fixed, and
\[
  Y:=n^\delta.
\]
Let $\mathcal P_Y$ be the set of odd primes at most $Y$.  For
$a<c$ in $\mathcal P_Y$, put
\[
  I_{a,c}:=\left(\frac{n}{2ac},\frac{n}{ac}\right],
  \qquad
  B_{a,c}:=I_{a,c}\cap\PP,
\]
and define the target family
\[
  \mathcal T
  :=
  \{\,acb:\ a<c,\ a,c\in\mathcal P_Y,\ b\in B_{a,c}\,\}.
\]
For large $n$, every such $b$ is greater than $Y$: indeed
$b>n/(2Y^2)=\frac12 n^{1-2\delta}>Y$.  Thus $a,c,b$ are distinct primes and
each target $acb$ lies in $(n/2,n]$.

\paragraph{Diagnostic auxiliary slot game and potential.}
The rank-three slot hypergraph $G(n,\delta)$ has vertex set consisting
of the slots $b,ab,cb$ which occur in the targets above, and has weighted edges
\[
  e_{a,c,b}:=\{b,ab,cb\},
  \qquad
  a<c,\quad a,c\in\mathcal P_Y,\quad b\in B_{a,c},
\]
each of weight $1$.  More generally, the finite auxiliary games below are
weighted incidence games of rank at most $3$.  A state has a captured-vertex set
$C$, a deleted-vertex set $D_V$, an unscored deleted-edge set $D_E$, a scored
edge set $K$, and total scored weight $S$.  An edge is live if it is not in
$D_E\cup K$ and contains no deleted vertex.  Maker may capture a live edge,
adding its vertices to $C$ and adding its weight to $S$; in the scored
hypergraph variant, Maker may instead use an \emph{alternate-scoring move} to
score the target without adding slot vertices to $C$.  Breaker may reply by
deleting one uncaptured vertex, by deleting one live edge without scoring it, or
in the scored-hypergraph variant by making a scored-edge reply, which adds one
live edge to $K$ and increments $S$ by its weight.  The scaled potential is
\[
  Q:=8S+\sum_{e\text{ live}}2^{|e\cap C|}w(e).
\]
The option of deleting an exact edge without adding its weight to $S$ is an
artificial strengthening.  It is useful for testing potential arguments, but
it is not a move of the residual divisibility game: playing an exact target
adds one to the actual game length.

\begin{proposition}[Max-gain domination is false in rank three]
\label{prop:t2-max-gain-counterexample}
The following finite scored rank-$3$ capture state has a unique max-gain Maker
move, but Breaker can answer that move by deleting one uncaptured vertex and
make the scaled potential decrease.
\end{proposition}

\begin{proof}
Let the already-captured vertices be
\[
  C=\{z_{ij}:1\le i,j\le 3\},
\]
and let $u_1,u_2,u_3,v$ be uncaptured.  There are no deleted or scored edges.
The live edges are
\[
  f=\{u_1,u_2,u_3\},\qquad w(f)=9,
\]
and, for $1\le i,j\le 3$,
\[
  e_{ij}=\{v,u_i,z_{ij}\},\qquad w(e_{ij})=4.
\]
For a live edge $e$, write $\Phi(e)=2^{|e\cap C|}w(e)$ and
$Q=8S+\sum\Phi(e)$.

Before Maker moves, $\Phi(f)=9$ and $\Phi(e_{ij})=8$ for all $i,j$.  If Maker
plays $f$, then the score contribution of $f$ changes by
$8w(f)-\Phi(f)=72-9=63$, and each of the nine edges $e_{ij}$ doubles from
$8$ to $16$.  Hence
\[
  \Delta(f)=63+9\cdot 8=135.
\]
After this move, the vertex $v$ is still uncaptured, and deleting $v$ removes
all nine edges $e_{ij}$, now of potential $16$ each.  The Breaker loss is
\[
  9\cdot 16=144>135.
\]
Thus the full round decreases $Q$.

It remains only to check that $f$ is the unique max-gain move.  If Maker plays
some $e_{ij}$, its own score gain is $8\cdot4-8=24$.  The edge $f$ doubles
from $9$ to $18$, giving gain $9$.  The two other edges with the same $u_i$
gain $24$ each, and the six edges with the other two $u$-vertices gain $8$
each.  Therefore
\[
  \Delta(e_{ij})=24+9+2\cdot24+6\cdot8=129<135.
\]
So the unique max-gain move is $f$, and it does not dominate the subsequent
vertex deletion at $v$.
\end{proof}

\begin{proposition}[\texorpdfstring{$K_4$}{K4}-fiber refutation of unscored deletion]
\label{prop:t2-safe-edge-refutation}
For every fixed $0<\delta<1/4$ and all sufficiently large $n$, the residual
slot hypergraph $G(n,\delta)$ supports a state reachable in the diagnostic
auxiliary game at which every legal Maker move admits an unscored edge-deletion
reply that strictly decreases the scaled potential $Q$.  Hence the general
safe-edge hypothesis is false for that stronger auxiliary game.
\end{proposition}

\begin{proof}
Here ``reachable'' means reachable under some legal play sequence, not
an actual divisibility-game history: the construction uses the auxiliary
capture, alternate-scoring, and unscored-deletion rules above.

Choose the four fixed small primes $13,17,19,23$.  Their pairwise products are
\[
  13\cdot17=221,\quad 13\cdot19=247,\quad 13\cdot23=299,\quad
  17\cdot19=323,\quad 17\cdot23=391,\quad 19\cdot23=437.
\]
All six products lie between $221$ and $437$, and $2\cdot221=442>437$.  For
sufficiently large $n$, since $23\le n^\delta=Y$, all four primes lie in
$\mathcal P_Y$.  By the prime number theorem, the interval
\[
  \left(\frac{n}{442},\frac{n}{437}\right]
\]
contains a prime $q$ for all sufficiently large $n$.  Also $q>Y$ for large
$n$, since $\delta<1/4$.

For every pair $a<c$ among $\{13,17,19,23\}$, we have
\[
  q>\frac n{442}\ge\frac n{2ac}
  \qquad\text{and}\qquad
  q\le\frac n{437}\le\frac n{ac},
\]
because $ac\ge221$ and $ac\le437$.  Hence $q\in B_{a,c}$ for all six pairs,
and the six edges
\[
  e_{a,c}:=\{q,aq,cq\}
\]
all occur in $G(n,\delta)$.  They form a $K_4$-fiber over the common slot $q$.

We next reach a state in which only these six edges are live and no vertex in
them has yet been captured.  While any other edge is live, Maker uses the
alternate-scoring move to score that exact target without adding slot vertices
to $C$.  Any edge outside this six-edge family has at least one slot outside
\[
  U_{\mathrm{fib}}:=\{q,13q,17q,19q,23q\},
\]
so Breaker deletes such an outside slot.  This never harms the six desired
edges.  Repeating this finite operation leaves exactly the six $K_4$-fiber
edges live, with $C=\varnothing$.

From that state, play two legal rounds.  First Maker captures
\[
  e_{13,17}=\{q,13q,17q\},
\]
so $q,13q,17q\in C$, and Breaker deletes the live edge $e_{13,19}$.  Next
Maker captures
\[
  e_{19,23}=\{q,19q,23q\},
\]
so $q,13q,17q,19q,23q\in C$, and Breaker deletes the live edge $e_{13,23}$.

At the next pre-Maker state, the only remaining live edges from this fiber are
\[
  e_{17,19}=\{q,17q,19q\},
  \qquad
  e_{17,23}=\{q,17q,23q\}.
\]
Both have all three vertices in $C$, so $\Phi(e)=2^3\cdot1=8$ for each.  There
are no other live edges by construction.

Maker has only two possible live edges to play.  If Maker plays $e_{17,19}$,
by ordinary capture or by alternate scoring, then the change in $Q$ is
\[
  +8\cdot1-8=0,
\]
because $S$ increases by $1$, the edge of potential $8$ is removed, and no new
vertex is added to $C$.  The remaining live edge $e_{17,23}$ still contributes
$8$.  Breaker then deletes $e_{17,23}$, decreasing $Q$ by $8$.  Thus
$Q_{\mathrm{after}}=Q_{\mathrm{before}}-8$.  The same argument holds if Maker
instead plays $e_{17,23}$.

Therefore this reachable pre-Maker state has the claimed property: every legal
Maker move admits a legal Breaker reply that strictly decreases $Q$.
Note that Breaker's final reply here is an unscored edge deletion, adding
$e_{17,23}$ to $D_E$, not a scored-edge play; under the scored-edge rule, the
same move would change $Q$ by $0$, not by $-8$.  The refutation shows that the
general safe-edge hypothesis fails when Breaker has access to the artificial
unscored-edge-deletion move.  In the actual residual game this reply is the
exact target $17\cdot23\cdot q$; it contributes one game move, so the
exact-move potential used below changes by $0$, not negatively.
\end{proof}

\begin{proposition}[Weighted activation max-threat selector]
\label{prop:t2-activation-selector}
Let a finite simple graph have positive weights $w_e$ on its live edges and a
shadow-captured vertex set $C$.  Put
\[
  \phi(e):=2^{|e\cap C|-3}w_e
\]
for each live edge, and, for every $v\notin C$, put
\[
  h(v):=\sum_{\substack{e\text{ live}\\v\in e}}\phi(e),\qquad
  H:=\max_{v\notin C}h(v),
  \qquad
  K:=\max_{e\text{ live}}\phi(e),
\]
with $H=0$ if every vertex is captured.  If at least one live edge exists,
there is a live edge $f$ such that the bank gain from claiming $f$, capturing
its endpoints, and replacing its contribution $\phi(f)$ by $w_f$ dominates
both of the following post-move threats:
\begin{enumerate}
  \item the entire live star at any still-uncaptured vertex;
  \item the coefficient of any remaining live edge.
\end{enumerate}
Consequently the bank cannot decrease if the reply deletes one uncaptured
vertex or one live edge.
\end{proposition}

\begin{proof}
Suppose first that $H\ge K$.  Choose $v\notin C$ with $h(v)=H$, and choose
an incident live edge $f=uv$ whose coefficient is maximal among the edges
incident with $v$.

If $u\notin C$, then $f$ has coefficient $w_f/8$.  Claiming it captures both
endpoints, so the exact bank gain before the reply is
\[
  G=H+h(u)+\frac58w_f.
\]
Indeed, the chosen contribution rises from $w_f/8$ to $w_f$, while the other
edges incident with $u$ or $v$ double.  For any still-uncaptured vertex $z$,
its post-move star is at most
\[
  h(z)+\phi(zv)+\phi(zu)
  \le H+\phi(f)+h(u)
  \le G.
\]
Here absent edges contribute zero.  A remaining edge through $v$ has
post-move coefficient at most $2\phi(f)=w_f/4$.  A remaining edge through
$u$ has coefficient at most $2h(u)$, and
\[
  2h(u)\le H+h(u)\le G.
\]
Every edge through neither new endpoint remains at most $K\le H\le G$.

If $u\in C$, then $\phi(f)=w_f/4$, only $v$ is newly captured, and the exact
gain is
\[
  G=H+\frac12w_f.
\]
A still-uncaptured star is at most
$H+\phi(f)=H+w_f/4$, a remaining edge through $v$ has coefficient at most
$2\phi(f)=w_f/2$, and every other edge remains at most $K\le H$.  Again all
post-move threats are at most $G$.

Finally suppose that $H<K$.  Any edge attaining $K$ has both endpoints in
$C$, since an uncaptured endpoint would have star at least $K$.  Claim such an
edge.  Its coefficient is $w_f/2$, no other coefficient changes, and replacing
$w_f/2$ by $w_f$ gains exactly $K$.  Every remaining vertex or edge threat is
at most $\max\{H,K\}=K$, proving the claim.
\end{proof}

\begin{proposition}[Initial target mass]\label{prop:t2-initial-mass}
For every fixed $0<\delta<1/4$,
\[
  W_0
  :=
  \sum_{\substack{a<c\\a,c\in\mathcal P_Y}} |B_{a,c}|
  \gg_\delta
  \frac{n(\log\log n)^2}{\log n}.
\]
\end{proposition}

\begin{proof}
Since $ac\le Y^2=n^{2\delta}$, the interval endpoint
$X_{a,c}:=n/(ac)$ satisfies $X_{a,c}\ge n^{1-2\delta}\to\infty$.  By the prime
number theorem, uniformly for all such pairs and all sufficiently large $n$,
\[
  \pi(X_{a,c})-\pi(X_{a,c}/2)
  \ge
  \kappa_\delta \frac{X_{a,c}}{\log n}
  =
  \kappa_\delta \frac{n}{ac\log n}
\]
for some positive constant $\kappa_\delta$.  Therefore
\[
  W_0
  \ge
  \kappa_\delta\frac{n}{\log n}
  \sum_{\substack{a<c\\a,c\in\mathcal P_Y}}\frac1{ac}.
\]
Mertens' theorem for primes gives
\[
  \sum_{p\in\mathcal P_Y}\frac1p
  =
  \log\log Y+O(1)
  =
  \log\log n+O_\delta(1),
\]
and $\sum_p p^{-2}<\infty$.  Hence
\[
  \sum_{a<c\in\mathcal P_Y}\frac1{ac}
  =
  \frac12\left[
    \left(\sum_{p\in\mathcal P_Y}\frac1p\right)^2
    -
    \sum_{p\in\mathcal P_Y}\frac1{p^2}
  \right]
  =
  \left(\frac12+o(1)\right)(\log\log n)^2.
\]
Combining the last two displays proves the claim.
\end{proof}

\begin{proposition}[Residual divisibility embedding]\label{prop:t2-residual-embedding}
Let $\mathcal T_*$ be a subfamily of targets $t=acb\in\mathcal T$ with the
following secured-pair origin.  For every $t=acb\in\mathcal T_*$, there is a
securing target
\[
  t_0(a,c)=acb_0\in\mathcal T,\qquad b_0\ne b,
\]
which has already been played, making $a$, $c$, and $ac$ unavailable as move
options; and no other earlier move is comparable with any
$t\in\mathcal T_*$.  Attach to $t$ the three slot labels
\[
  b,\qquad ab,\qquad cb
\]
and let $e_t:=\{b,ab,cb\}$.  Then the residual divisibility game on
$\mathcal T_*$ is at least as favorable to Prolonger as the scored rank-$3$
hypergraph game on the hyperedges $\{e_t:t\in\mathcal T_*\}$:
\begin{enumerate}
  \item the only future moves that can make $t=acb$ illegal are
    $b$, $ab$, $cb$, and $t$ itself;
  \item playing $b$, $ab$, or $cb$ deletes exactly the residual targets whose
    hyperedges contain that slot;
  \item playing the exact target $t$ kills no other target in $\mathcal T_*$,
    and is therefore a scored-edge move;
  \item every live hyperedge corresponds to a legal actual move $t$.
\end{enumerate}
\end{proposition}

\begin{proof}
Fix $t=acb\in\mathcal T_*$.  Because $t>n/2$, no proper multiple of $t$ lies in
$\{2,\dots,n\}$.  Thus a future move can make $t$ illegal only by being $t$
itself or a proper divisor of $t$.  The positive divisors are
\[
  1,\ a,\ c,\ ac,\ b,\ ab,\ cb,\ acb.
\]
The divisor $1$ is not a move, while $a$, $c$, and $ac$ are unavailable by
the played securing target $t_0(a,c)$.  This proves the first assertion.

For slot incidence, suppose another target is $t'=a'c'b'$.  Since
$b,b'>Y\ge a,a',c,c'$, the divisibility $b\mid t'$ holds iff $b=b'$.  Thus a
move $b$ kills exactly the targets in the $b$-fiber.  Similarly, if
$ab\mid a'c'b'$, the large prime factor forces $b=b'$, and after cancellation
$a$ must be $a'$ or $c'$.  This is exactly the condition that the slot $ab$
belongs to $e_{t'}$.  The proof for $cb$ is identical.

If two targets $t,u$ in $\mathcal T_*$ are comparable, say $t\mid u$, then
$u$ is a multiple of $t$ lying in $(n/2,n]$.  The next positive multiple of
$t$ after $t$ is $2t>n$, so $u=t$.  Hence distinct targets are pairwise
incomparable, and playing an exact target scores only that target.

Finally, if $e_t$ is live in the slot hypergraph, then none of $b$, $ab$, $cb$,
or $t$ has been played.  The only other proper divisors of $t$ are
$1,a,c,ac$, and these are either not moves or are unavailable because of the
played securing target $t_0(a,c)$.  The securing target is a distinct target
larger than $n/2$, hence is incomparable with $t$, and there is no proper
multiple of $t$ on the board.  By hypothesis, no other earlier move is
comparable with any target in $\mathcal T_*$.  Therefore no played number is
comparable with $t$, so the actual move $t$ is legal.  This proves the
comparison.
\end{proof}

\begin{proposition}[Activation-stage potential]\label{prop:t2-activation-potential}
Throughout, a target is \emph{live} if no previously played integer is
comparable with it under divisibility; a token is deleted (its target is
\emph{killed}) precisely when its target ceases to be live.
During activation, let each unclaimed pair edge $e=(a,c)$ carry the current
token set $B_e(t)$ of primes $b\in B_{a,c}$ for which the target $acb$ has not
yet been killed.  Write $w_t(e):=|B_e(t)|$; after $e$ is secured, use the same
notation for its remaining live residual tokens.

Prolonger plays as follows.  In each activation round, among pair edges with
positive current token weight, he chooses the edge supplied by
\Cref{prop:t2-activation-selector}; if $e=(a,c)$ is chosen, he plays one live
target $acb$ with $b\in B_e(t)$.  This claims the pair edge $e$ and leaves the
other live tokens on that secured pair for the residual phase.

Let $C_t$ be the conservative shadow-capture set consisting only of small-prime
endpoints of Prolonger's securing targets; endpoint unavailability caused by
Shortener is ignored.  Let $T_t$ be the number of actual game moves made by
either player since activation began.  Define
\[
  \Psi_t
  :=
  T_t
  +\sum_{e\text{ secured}}w_t(e)
  +\sum_{e\text{ live}}\phi_t(e),
\]
where
\[
  \phi_t(e):=2^{|e\cap C_t|-3}w_t(e)
\]
for a live unclaimed pair edge.
Let $E$ be the total number of target tokens deleted during activation by
Shortener moves not modeled as graph vertex deletions $a$ or graph edge
deletions $ac$.  Here $E$ counts the total number of live target tokens deleted
by off-model Shortener moves during activation, not the number of such moves:
a single off-model move can delete multiple tokens, and $E$ tallies deletions,
not move events.  The response to the final Prolonger activation move is
included before the residual mass is frozen: the residual family is frozen
only after Shortener's reply to the final activation move, so the frozen
residual state is at a Prolonger turn.  If Shortener has no legal reply, the
whole game has ended, and then no residual target is live, since a live
residual target would itself be a legal reply.  At the end of activation,
\[
  T_{\mathrm{act}}
  +
  \sum_{e\text{ secured}} w_{\mathrm{end}}(e)
  \ge
  \frac{W_0}{8}-E.
\]
\end{proposition}

\begin{proof}
At each activation state, form the finite rank-$2$ graph game whose vertices are
the small primes, whose live edges are the unclaimed pairs $e=(a,c)$ with
current weight $w_t(e)>0$.  The captured vertices in the selector are exactly
$C_t$.  Shortener-created endpoint unavailability is favorable to Prolonger;
ignoring it only gives the model more future deletion options.  A modeled
Shortener reply of type $a$ or $c$ is a graph vertex deletion, and a modeled
reply of type $ac$ is a graph edge deletion.

Initially $\Psi_0=W_0/8$.  First ignore off-model token deletions.  When
Prolonger claims an edge $e$ of current weight $w_t(e)$, his exact target adds
one to $T_t$ and the other $w_t(e)-1$ tokens become secured.  Thus the chosen
edge's total bank contribution changes from $\phi_t(e)$ to
\[
  1+(w_t(e)-1)=w_t(e),
\]
while all other incident coefficients change exactly as in
\Cref{prop:t2-activation-selector}.  That proposition makes the Prolonger gain
dominate every modeled vertex or edge deletion.  Shortener's reply also adds
one actual move to $T_t$, which can only improve the inequality.  Hence
$\Psi_t$ is nondecreasing through every modeled full round.

Now restore all actual Shortener moves.  A token $acb$ is deleted precisely
when Shortener plays one of its still-available harmful divisors or the target
itself.  Deleting one live token can reduce $\Psi_t$ by at most $1$, since a
secured token has coefficient $1$ and every live-edge coefficient is at most
$1/2$.
Therefore all off-model deletions together reduce the lower bound by at most
$E$, giving
\[
  \Psi_{\mathrm{end}}\ge W_0/8-E.
\]
Each activation round is frozen only after Shortener's reply, if a reply
exists; in particular, a last reply that exhausts the unclaimed graph is
already included in $T_{\mathrm{act}}$ and $E$.  At the end no positive-weight
unclaimed edge remains, so $\Psi_{\mathrm{end}}$ is exactly
$T_{\mathrm{act}}$ plus the residual live target weight on secured pairs.
\end{proof}

\begin{proposition}[Activation deletion budget]\label{prop:t2-deletion-budget}
For fixed $0<\delta<1/4$, the off-model deletion count from
\Cref{prop:t2-activation-potential} satisfies
\[
  E\ll \frac{Y^4}{\log^4 Y}
  =
  o\!\left(\frac{n(\log\log n)^2}{\log n}\right).
\]
\end{proposition}

\begin{proof}
There is at most one activation round per small-prime pair, so the number of
rounds is
\[
  R\le \binom{|\mathcal P_Y|}{2}
  \ll \frac{Y^2}{\log^2Y}.
\]
We bound the number of live tokens an off-model Shortener move can delete.

If Shortener plays a large prime $b$, then for each pair $(a,c)$ there is at
most one token $acb$ using that $b$.  Hence the move deletes at most
$O(Y^2/\log^2Y)$ tokens.  If Shortener plays a lateral semiprime $pb$ with
$p\le Y$ and $b>Y$ prime, the killed targets have the form $pcb$ with the
third small prime varying, so there are at most $O(Y/\log Y)$ of them.  If
Shortener plays an exact target $acb$, exactly one token is deleted.

The worst of these bounds is the large-prime bound.  Across $R$ activation
rounds,
\[
  E\ll R\frac{Y^2}{\log^2Y}
  \ll \frac{Y^4}{\log^4Y}.
\]
Since $Y=n^\delta$,
\[
  \frac{Y^4/\log^4Y}{n(\log\log n)^2/\log n}
  \ll_\delta
  \frac{n^{4\delta-1}}{(\log n)^3(\log\log n)^2}
  \to 0
\]
because $4\delta-1<0$.
\end{proof}

\begin{proposition}[Residual mass on secured pairs]\label{prop:t2-residual-mass}
At the end of activation, the secured pairs carry residual live target weight
\[
  M:=
  \sum_{e\text{ secured}} w_{\mathrm{end}}(e)
  \gg_\delta
  \frac{n(\log\log n)^2}{\log n}.
\]
\end{proposition}

\begin{proof}
By \Cref{prop:t2-activation-potential},
\[
  T_{\mathrm{act}}+M\ge \frac{W_0}{8}-E.
\]
There is at most one Prolonger activation move per small-prime pair and at
most one Shortener reply to each such move, so
\[
  T_{\mathrm{act}}\le 2R\ll \frac{Y^2}{\log^2Y}
  =
  o\!\left(\frac{n(\log\log n)^2}{\log n}\right),
\]
because $2\delta<1$.  The deletion budget $E$ is negligible by
\Cref{prop:t2-deletion-budget}, while $W_0$ has the required lower bound by
\Cref{prop:t2-initial-mass}.  Subtracting the two negligible terms gives the
claim.
\end{proof}

\begin{proposition}[Exact-move residual fiber capture]
\label{prop:t2-residual-capture}\label{prop:t2-finite-capture}
Let $\mathcal T_*$ satisfy the secured-pair hypotheses of
\Cref{prop:t2-residual-embedding}, suppose $|\mathcal T_*|=M$, and suppose the
frozen residual state is at a Prolonger turn.  From this state Prolonger has a
strategy under which at least $M/8$ exact targets in $\mathcal T_*$ are played
by one of the two players before no target in $\mathcal T_*$ remains live.
\end{proposition}

\begin{proof}
For each large prime $q$, form the simple side graph $G_q$: its vertices are
the side slots $aq$ which occur in $\mathcal T_*$, and an edge $ac$ represents
the target $acq$.  Distinct $q$-fibers have disjoint side slots by unique
factorization.  At the residual phase boundary we reset the shadow state:
every surviving fiber is declared unactivated and every side-capture set is
empty.  Thereafter the shadow records only Prolonger targets played from
$\mathcal T_*$.  Additional slot unavailability created by activation-stage
moves or incidentally by a Shortener exact-target play is deliberately ignored;
this only gives the modeled Shortener more reply options.  A direct Shortener
common- or side-slot play still deletes its incident current edges, as checked
below.

Before Prolonger has subsequently played a target in $G_q$, call the fiber
unactivated and
give every live edge coefficient $1/8$.  After the first such play, call it
activated: the common slot $q$ is captured, and a live side edge $e$ has
coefficient
\[
  2^{c_q(e)-2}\in\left\{\frac14,\frac12,1\right\},
\]
where $c_q(e)$ is the number of its two side endpoints shadow-captured by
Prolonger.  At each time let $E_q$ denote the current set of live edges in
$G_q$.  Let $T_{\mathrm{res}}$ count every exact target in
$\mathcal T_*$ played by either player, and define
\[
  \Psi_{\mathrm{res}}
  :=
  T_{\mathrm{res}}
  +\frac18\sum_{q\text{ unactivated}}|E_q|
  +\sum_{\substack{q\text{ activated}\\e\in E_q}}
     2^{c_q(e)-2}.
\]
Initially $\Psi_{\mathrm{res}}=M/8$.

For an uncaptured side vertex $v$ in an activated fiber, let
\[
  g_q(v):=
  \sum_{\substack{e\in E_q\\v\in e}}2^{c_q(e)-2}.
\]
At each Prolonger turn, compare all activated threats $g_q(v)$ with all
unactivated common-slot threats $|E_q|/8$, and let $H$ be their maximum, with
$H:=0$ if this collection is empty.

Suppose first that $H=g_q(v)>0$ in an activated fiber.  If $v$ has an
uncaptured neighbor $u$, Prolonger plays the target represented by $uv$.  The
edge has coefficient $1/4$, both endpoints become captured, and the exact
increase before Shortener's reply is
\[
  g_q(v)+g_q(u)+\frac14\ge H+\frac12.
\]
The threat at any other uncaptured vertex can rise only through its possible
edges to $u$ and $v$, by at most $1/4$ for each, so every post-move side threat
is at most $H+1/2$.  If $v$ has no uncaptured neighbor, Prolonger instead plays
an incident edge whose other endpoint is captured.  Its coefficient is
$1/2$, and the exact gain is
\[
  1-\frac12+\left(g_q(v)-\frac12\right)=g_q(v)=H.
\]
No other uncaptured side threat rises in this case.

Now suppose that $H=m/8>0$ is attained by an unactivated fiber with $m$ live
edges.  Choose a maximum-degree side vertex $u$ and an incident edge $uv$;
write their degrees as $d_u,d_v$.  Activating this fiber changes its common
coefficient from $1/8$ per edge to the activated coefficients and gives the
exact gain
\[
  G=\frac m8+\frac{d_u+d_v}{4}+\frac14.
\]
Every new uncaptured side threat is at most $d_u/4+1/2$, and
\[
  \frac{d_u}{4}+\frac12
  \le
  \frac m8+\frac{d_u+d_v}{4}+\frac14=G,
\]
because $m,d_v\ge1$.  All threats in other fibers were at most
$H=m/8\le G$.

It remains to check that these are all harmful Shortener replies.  In an
unactivated fiber, playing the common slot $q$ removes exactly
$|E_q|/8$ from the bank, while playing a side slot removes at most that
fiber's side-star contribution.  In an activated fiber the common slot is
already unavailable, and a legal side-slot reply removes precisely a
post-move threat bounded above.  By
\Cref{prop:t2-residual-embedding}, the only other harmful reply is an exact
target.  Such a reply removes an edge of coefficient at most $1$ and
simultaneously increases $T_{\mathrm{res}}$ by $1$, so it never lowers
$\Psi_{\mathrm{res}}$.  Replies unrelated to $\mathcal T_*$ do not change the
bank.  Thus the chosen Prolonger gain covers every possible decrease during
the following reply.

If $H=0$ while a live edge remains, every fiber is activated and both
endpoints of every live edge are shadow-captured.  Prolonger plays any such
edge; its coefficient $1$ is replaced by the unit increase in
$T_{\mathrm{res}}$, leaving the bank unchanged.  Its common and side slots
are actual divisors of earlier Prolonger targets and hence unavailable.

Therefore $\Psi_{\mathrm{res}}$ is nondecreasing until no residual edge
remains.  A live edge always corresponds to a legal actual target by
\Cref{prop:t2-residual-embedding}, so the actual game cannot terminate sooner.
At exhaustion,
\[
  T_{\mathrm{res}}
  =\Psi_{\mathrm{res,fin}}
  \ge\Psi_{\mathrm{res},0}
  =\frac M8,
\]
as required.
\end{proof}

\begin{proposition}[Unconditional rank-three lower bound]\label{prop:t2-final}
For every fixed $0<\delta<1/4$ there is $c_\delta>0$ such that
\[
  \L(n)\ge c_\delta\frac{n(\log\log n)^2}{\log n}
\]
for all sufficiently large $n$.
\end{proposition}

\begin{proof}
After activation, restrict to the residual live targets lying over secured
pairs.  Their total weight is $M$, as in
\Cref{prop:t2-residual-mass}.  For each secured pair $(a,c)$, the activation
move supplies a securing target $t_0(a,c)=acb_0$, and the residual subfamily
keeps only targets $acb$ with $b\ne b_0$ which remain live.  Thus the securing
target makes $a$, $c$, and $ac$ unavailable while being incomparable with every
remaining residual target over the same pair.  Any activation move comparable
with a target removes that target before $\mathcal T_*$ is frozen.  Thus
$\mathcal T_*$ satisfies the hypotheses of
\Cref{prop:t2-residual-embedding}, and its cardinality is exactly $M$: since
$b>Y\ge a,c$, unique factorization shows each target $acb$ determines the
triple $(a,c,b)$, so distinct live tokens over secured pairs represent
distinct targets and $|\mathcal T_*|=\sum_{e\text{ secured}}
w_{\mathrm{end}}(e)=M$.

By \Cref{prop:t2-residual-capture}, at least $M/8$ exact targets from this
subfamily are then played by Prolonger or Shortener.  Each is a genuine actual
game move.  Hence the total game length is at least $M/8$.  Since
$M\gg_\delta n(\log\log n)^2/\log n$, the same lower bound holds for the
total game length unconditionally.

This proves the per-$\delta$ statement \Cref{thm:t2}; the intro's absolute-$c$
corollary follows by specializing $\delta=1/8$.
\end{proof}


\subsection{Diagonal selection conditions}\label{app:diagonal-selection-conditions}

This subsection records the eventual side conditions used in
\Cref{prop:prime-bridge-diagonalization}.  In the notation of that proposition,
after $H_m,\lambda_m,a_m,\eta_m,s_m$ have been fixed, choose $N_m$ recursively
increasing so that for every $n\ge N_m$ the following twelve conditions hold:
\begin{enumerate}
  \item the prefix-existence conclusion of \Cref{lem:prefix-existence} holds
  up to index $K$;
  \item the local-density error from \Cref{prop:local-density} satisfies
  \[
    \xi_{H_m}(n)\le \frac1m;
    \tag{7.10a}\label{eq:diag-xi}
  \]
  \item the envelope moment estimate from \Cref{prop:envelope-inversion}
  satisfies
  \[
    |T_r^{(b^{(H_m)})}(n)-J_r^{(H_m)}|\le \frac{1}{4m}
    \qquad (1\le r\le4);
    \tag{7.10b}\label{eq:diag-envelope}
  \]
  \item the bridge moment estimate from \Cref{prop:queued-prime-rounding}
  satisfies
  \[
    |T_r^{(p^{(H_m)})}(n)-T_r^{(b^{(H_m)})}(n)|\le \frac{1}{2m}
    \qquad (1\le r\le4);
    \tag{7.10c}\label{eq:diag-bridge}
  \]
  \item the cell-reciprocal-mass estimate of
  \Cref{lem:cell-reciprocal-mass} holds for the finitely many fixed-ratio
  cells used with parameters $H_m,a_m,s_m$, with all limiting errors bounded
  by $1/m$;
  \item the product-strip estimate of \Cref{lem:product-strip-reciprocal}
  holds for $C=\lambda_m^r$ and $1\le r\le4$, with associated finite
  cell-count estimates valid to tolerance $1/m$;
  \item the queue-clearing Hall inequalities in
  \Cref{prop:queued-prime-rounding}, with $a=a_m$, $\eta=\eta_m$, and
  $s=s_m$, hold for every suffix $0\le t\le s_m$;
  \item the exceptional-index reciprocal mass estimate in
  \eqref{eq:exceptional-b-mass} holds in the form
  \[
    \sum_{j\in E}\frac1{b_j}\le \frac1m;
  \]
  \item all fixed-ratio prime number theorem estimates used in the queue
  construction, in particular \eqref{eq:prime-bin-count}, hold with tolerance
  $1/m$ for the finitely many ratios determined by $a_m$ and $s_m$;
  \item the smallest genuine scale satisfies
  \[
    n^{\alpha_{H_m-1}}
    =
    n^{1/H_m+\tau_{H_m}}
    \ge m;
    \tag{7.10d}\label{eq:diag-scale}
  \]
  \item the interblock nonoverlap inequalities
  \[
    a_m^2B_h<A_{h-1}\qquad(2\le h\le H_m-1),
    \qquad
    a_m^2B_1<n
  \]
  hold for the block endpoints associated to $H_m$;
  \item the top-block prime supply satisfies
  \[
    \pi(2n)-\pi(n)\ge K.
  \]
\end{enumerate}
Each condition is eventual because $H_m,\lambda_m,a_m,\eta_m,s_m$ and the
finite list of ratios are fixed while $n\to\infty$.


\subsection{Restricted carrier classes}\label{app:restricted-carriers}

The next two propositions record unconditional upper bounds for two restricted
families of Prolonger carrier moves.  They are not used in the main upper
bound, but they clarify which carrier configurations are not responsible for
the remaining difficulty.

\begin{proposition}[Small-prime composite Prolonger with disjoint supports]
\label{prop:disjoint-carriers}
Fix $1/3<\alpha<1/2$ and put $y=n^\alpha$.  Suppose that every composite
Prolonger move has all prime factors at most $y$, and that the prime supports
of distinct composite Prolonger moves are pairwise disjoint.  Against this
restricted class, Shortener has a strategy forcing
\[
  \L(n)=O_\alpha\!\left(\frac n{\log n}\right).
\]
\end{proposition}

\begin{proof}
Shortener uses the following priority list, always skipping moves which have
already become illegal.  First, she plays every legal prime.  After this prime
phase, define $B$ to be the set of primes $p\le y$ whose singleton was not
resolved by prime play because $p$ appears in an observed composite Prolonger
carrier.  By the disjoint-support hypothesis, each $p\in B$ belongs to a
unique composite carrier $C(p)$.

Second, for every $p\in B$ whose carrier has at least two distinct prime
factors, write $e(p)=v_p(C(p))$.  If $p^{e(p)+1}\le n$, Shortener plays
$p^{e(p)+1}$.  This move is legal when it is attempted: it is not comparable
with $C(p)$, because it has a larger $p$-adic exponent while $C(p)$ contains
another prime factor; it is not comparable with any other composite carrier by
disjointness of supports; and previous priority moves either contain different
prime supports or have smaller rank.

Third, for every pair $p,q\in B$ belonging to different carriers, Shortener
plays the legal squarefree semiprime $pq$ when possible.  Such a move is
incomparable with every carrier, since no carrier contains both $p$ and $q$,
and neither $p$ nor $q$ alone is available as a legal prime at this stage.

We use the following skipped-move invariant.  Under the disjoint-support
hypothesis, no prior composite Prolonger move is a proper multiple of a planned
Shortener move in the second or third phase.  For the phase-two move
$p^{e(p)+1}$, the unique carrier containing $p$ has only $p$-adic exponent
$e(p)$, while all other carriers omit $p$.  For the phase-three move $pq$, the
two primes lie in different carriers, so no single carrier contains both.  Thus
when a planned move is skipped because it is already illegal, the earlier
comparable move may be treated as a divisor of the planned move, or the planned
move itself, but not as a proper composite multiple coming from Prolonger.

The number of moves in the three phases is
\[
  O(\pi(n))+\pi(y)+\pi(y)^2
  =
  O_\alpha\!\left(\frac n{\log n}\right),
\]
because $2\alpha<1$.

We now show that after these phases every future legal move has no prime
factor at most $y$.  Let $x$ be a legal move divisible by a prime $p\le y$.

If $p\notin B$, then phase one either played $p$ or $p$ was already played by
Prolonger as a prime; in either case any multiple of $p$ is illegal.  Thus
$p\in B$.

If $x$ contains primes from two distinct carriers, then one of the phase-three
semiprimes $pq$ divides $x$, unless it was already illegal because an earlier
played move dividing $pq$ also divides $x$.  Either way $x$ is illegal.

Hence all prime factors of $x$ which are at most $y$ lie in a single carrier
support.  If the carrier is a pure prime power $p^e$, then any $x$ divisible by
$p$ and supported only on this carrier is another power of $p$, and is
comparable with $p^e$.  If the carrier has at least two distinct prime factors,
then either $x$ divides the carrier, in which case $x$ is illegal, or for some
$p$ it has exponent exceeding $v_p(C(p))$.  In the latter case
$p^{e(p)+1}\mid x$, so the phase-two move makes $x$ illegal.

Thus all primes $\le y$ are resolved.  Any remaining legal integer has all
prime factors greater than $y$.  Since $y^3>n$, it has at most two prime
factors.  The number of integers at most $n$ with one or two prime factors,
all exceeding $y=n^\alpha$, is $O_\alpha(n/\log n)$ by the prime number theorem
and the standard semiprime estimate
\[
  \sum_{p>y}\pi(n/p)\ll_\alpha \frac n{\log n}\sum_{p>y}\frac1p
  =O_\alpha\!\left(\frac n{\log n}\right),
\]
where the sum is over $p\le n/y$.  Adding the phase moves proves the claim.
\end{proof}

\begin{proposition}[Squarefree rank-three small-prime carriers]\label{prop:rank-three-carriers}
Fix $1/3<\alpha<1/2$ and put $y=n^\alpha$.  Suppose that every composite
Prolonger move is squarefree, has all prime factors at most $y$, and has
support size at most $3$.  Against this restricted class, Shortener has a
strategy forcing
\[
  \L(n)=O_\alpha\!\left(\frac n{\log n}\right).
\]
\end{proposition}

\begin{proof}
Let $B$ be the set of primes which occur in Prolonger's composite carriers.
Shortener again uses a fixed priority list, skipping illegal moves.

First she plays every legal prime.  Second she plays $p^2$ for every
$p\in B$ when legal.  Third she plays every legal squarefree semiprime $pq$
with $p,q\in B$.  Fourth she plays every legal squarefree triple $pqr$ with
$p,q,r\in B$.

The skipped-move invariant is as follows.  Under the squarefree rank-$\le3$
restriction, a prior Prolonger move which makes a planned square or planned
triple illegal cannot be a proper multiple of that planned move: a square has no
squarefree proper multiple, and a proper multiple of a squarefree triple has
rank at least $4$.  For a planned semiprime $pq$, the only possible proper
multiple is a rank-three carrier $pqr$; in the support-size-$2$ case below the
remaining candidate is exactly $pq$, hence it is already illegal because
$pq\mid pqr$.  Thus the later skipped-move arguments may charge a skipped
triple to a prior divisor or to the triple itself, and skipped semiprimes are
harmless in the only place where they are used.

The first phase has $O(n/\log n)$ moves, the second has at most $\pi(y)$ moves,
and the third has at most $\pi(y)^2=o(n/\log n)$ moves.  It remains to bound
the number of triples
\[
  T_\alpha(n):=\#\{p<q<r\le y:\ pqr\le n\}.
\]
Split according to whether $pq\le n^{1-\alpha}$.

If $pq\le n^{1-\alpha}$, then there are
$O(n^{1-\alpha}\log\log n/\log n)$ possible pairs $p<q$, and for each pair at
most $\pi(y)=O(n^\alpha/\log n)$ choices of $r$.  This contributes
$o(n/\log n)$.

If $pq>n^{1-\alpha}$, then $q\le y=n^\alpha$ implies
$p>n^{1-2\alpha}$.  In this range,
\[
  \#\{r:\ q<r\le y,\ pqr\le n\}
  \ll_\alpha \frac{n}{pq\log n}.
\]
Therefore the contribution is at most
\[
  \ll_\alpha
  \frac n{\log n}
  \sum_{n^{1-2\alpha}<p<q\le n^\alpha}\frac1{pq}
  =
  O_\alpha\!\left(\frac n{\log n}\right),
\]
because the reciprocal prime sum over a fixed-power interval is $O_\alpha(1)$.
Thus $T_\alpha(n)=O_\alpha(n/\log n)$.

After the prime phase, every prime $p\le y$ with $p\notin B$ has either been
played as the singleton $p$ or has had some proper multiple played; in either
case, no future legal Prolonger move can contain $p$ as a factor.  Therefore
any remaining composite Prolonger move is supported on $B$.

We prove that the priority list resolves all primes at most $y$.  After the
prime and square phases, any remaining legal move supported on $B$ must be
squarefree.  If it has support size $1$, it is comparable with a played or
blocked prime.  If it has support size $2$, then the corresponding semiprime
was either played in phase three or was already illegal by the skipped-move
invariant.  If it has support size $3$, the same statement holds for phase four,
where the invariant rules out proper multiples.

If a remaining legal move $x$ had support size at least $4$ inside $B$, choose
three primes $p,q,r$ from its support.  The triple $pqr$ was addressed in phase
four.  If it was played, then $pqr\mid x$ and $x$ is illegal.  If it was not
legal when considered, then some earlier played move comparable with $pqr$
must divide $pqr$ or equal $pqr$ by the skipped-move invariant; since all
earlier non-prime moves in the priority list have
support size at most $3$ and are squarefree except the squares $p^2$, that
played divisor also divides $x$.  Again $x$ is illegal.  Thus no remaining
legal move contains a prime at most $y$.

As in the proof of \Cref{prop:disjoint-carriers}, once all primes at most
$y$ are resolved, the residual legal set has size $O_\alpha(n/\log n)$.
Combining this residual bound with the phase counts proves the proposition.
\end{proof}
